\chapter{مقدمه}
\label{chap:intro}

نرم‌افزارهای دسکتاپ مبتنی بر
\lr{Qt C++}
سهم بزرگی از کاربردهای مهندسی، صنعتی و سازمانی را به خود اختصاص می‌دهند؛ از ابزارهای پایش شبکه گرفته تا نرم‌افزارهای کنترل تجهیزات. با این حال، بر خلاف برنامه‌های وب که اکوسیستم غنی‌ای از ابزارهای تست خودکار در اختیار دارند، تست خودکار رابط کاربری برای نرم‌افزارهای مبتنی بر
\lr{Qt}
تا حد زیادی به ابزارهای تجاری محدود مانده است. نتیجه‌ی این خلأ معمولاً تست دستی، پوشش بازگشت ناکافی و هزینه‌ی بالای نگهداری کیفیت در هر تغییر است.

هدف این پایان‌نامه، طراحی و پیاده‌سازی یک
\lr{\gls{E2E}}
چارچوب تست سرتاسری
برای نرم‌افزارهای مبتنی بر
\lr{Qt C++}
در بستر ویندوز است؛ چارچوبی که با ابزارهای متن‌باز ساخته شود، رفتار برنامه‌ی تحت تست را از دید کاربر واقعی بررسی کند و خروجی قابل‌فهمی برای تیم‌های توسعه و کیفیت تولید نماید. در این چارچوب، تعامل با برنامه از طریق
\gls{UiAutomation}
انجام می‌شود:
\lr{WinAppDriver}
درایور خودکارسازی ویندوز است و
\lr{Appium}
به عنوان سرور پروتکل بین تست‌ها و درایور عمل می‌کند
\cite{appium,winappdriver}.

میان‌بُرد اصلی این کار، ساختن یک
\lr{\gls{DSL}}
زبان مخصوص دامنه
در لایه‌ی بالای چارچوب است؛ به گونه‌ای که هر تست صرفاً با عباراتی مانند «واحد استان تهران را فیلتر کن، جستجو کن و تعداد ردیف‌ها را بخوان» نوشته شود و هیچ‌یک از تست‌ها مستقیماً با جزئیات درایور سروکار نداشته باشند. این جداسازی باعث می‌شود تغییرات رابط کاربری تنها در لایه‌ی اشیای صفحه
\gls{PageObject}
اعمال شود و سوییت‌ها دست‌نخورده بماند
\cite{fowler2002,meszaros2007}.

برای ارزیابی چارچوب، یک برنامه‌ی نمونه با عنوان «ردیاب سیگنال ایران» طراحی و پیاده‌سازی شد؛ یک سامانه‌ی پایش کیفیت سیگنال مخابراتی مبتنی بر
\lr{Qt Widgets}
و
\lr{SQLite}
با داشبورد، نمودارها، فیلترها و گزارش‌گیری
\cite{qtdocs,sqlite}. این برنامه به عنوان
\lr{\gls{SUT}}
بکار گرفته شده و مجموعه‌ای از ۲۵ سناریوی تست سرتاسری روی آن اجرا می‌شود که همه‌ی جریان‌های کاربر — از ورود و ثبت‌نام تا فیلترهای ترکیبی، صفحه‌بندی، نمودارها و گزارش‌گیری — را پوشش می‌دهند. خروجی هر اجرا در قالب گزارش
\lr{Allure}
منتشر می‌شود
\cite{allure}
و مسیر انتقال چارچوب به پایپلاین‌ی
\lr{\gls{CI}}
/ \lr{\gls{CD}}
نیز طراحی شده است.

ساختار ادامه‌ی پایان‌نامه چنین است: فصل
\autoref{chap:concepts}
مفاهیم پایه‌ای لازم را معرفی می‌کند؛ فصل
\autoref{chap:related}
مروری بر کارهای مشابه دارد؛ فصل
\autoref{chap:demoapp}
برنامه‌ی نمونه را معرفی می‌کند؛ فصل
\autoref{chap:framework}
طراحی و پیاده‌سازی چارچوب را تشریح می‌کند؛ فصل
\autoref{chap:evaluation}
ارزیابی و نتایج را ارائه می‌دهد و فصل
\autoref{chap:conclusion}
نتیجه‌گیری و کارهای آینده است.